\subsection{EventType}
\index{EventType: : : :!}The EventType class represents a type of event. Each
instance of an Event happening has one associated EventType, and callback
functions\index{callback function: : : :!} can be registered against
EventTypes. All EventTypes have an associate code$\text{\textgreek{—}}$an
integral value that identifies the EventType. Some EventTypes also have a
time associated with them (Pre, Post, or None)$\text{\textgreek{—}}$describing
when an Event may occur relative to the Code. For example, an EventType with
a code of Exit and a time of Pre (written as pre-exit) would be associated with
an Event that occurs just before a process exits and its address space is
cleaned. An EventType with code Exit and a time of Post would be associated
with an Event that occurs after the process exits and the address space is
cleaned.

When using EventTypes to register for callback functions a special time value of
Any can be used. This signifies that the callback function should trigger for
both Pre and Post time events. ProcControlAPI will never deliver an Event
that has an EventType with time code Any.

More details on Events and EventTypes can be found in Section
\ref{bkm:Ref273435580}.

EventType Types:

typedef enum \{

Pre = 0,

Post,

None,

Any

\} Time; \index{EventType, Time: : : :!}


\bigskip

typedef int Code;

The Time and Code types are respectively used to describe the time and code
values of an EventType. 

EventType Constants:

static const int Error = -1\index{EventType, Code: : : :!}

static const int Unset = 0

static const int Exit = 1

static const int Crash = 2

static const int Fork = 3

static const int Exec = 4

static const int ThreadCreate = 5

static const int ThreadDestroy = 6

static const int Stop = 7

static const int Signal = 8

static const int LibraryLoad = 9

static const int LibraryUnload = 10

static const int Bootstrap = 11

static const int Breakpoint = 12

static const int RPC = 13

static const int SingleStep = 14

static const int Library = 15

static const int MaxProcCtrlEvent = 1000

These constants describe possible values for an EventType$\text{\textgreek{’}}$s
code. The Error and Unset codes are for handling error cases and should not
be used for callback functions or be associated with Events. 

The EventType codes were implemented as an integer (rather than an enum) to
allow users to create custom EventTypes. Any custom EventType should begin at
the MaxProcCtrlEvent value, all smaller values are reserved by ProcControlAPI.

EventType Related Types:

struct eventtype\_cmp\index{eventtype\_cmp: : : :!} \{

bool operator()(const EventType \&a, const EventType \&b);

\}

This type defines a less-than comparison function for EventTypes. While a
comparison of EventTypes does not have a semantic meaning, this can be useful
for inserting EventTypes into maps or other STL data structures.

EventType Member Functions:

EventType\index{EventType constructor: : : :!}(Code e)

Constructs an EventType with the given code and a time of Any.

EventType\index{EventType constructor: : : :!}(Time t, Code e)

Constructs an EventType with the given time and code values.

EventType\index{EventType constructor: : : :!}()

Constructs an EventType with an Unset code and None time value.

Code code\index{code, EventType: : : :!}() const

Returns the code value of the EventType.

Time time\index{time, EventType: : : :!}() const

Returns the time value of the EventType.

std::string name\index{name, EventType: : : :!}() const

Returns a human readable name for this EventType.

